\documentclass[11pt]{article}

\usepackage[a4paper,margin=1in]{geometry}
\usepackage{amsmath,amssymb}
\usepackage{microtype}

\title{Discussion Report on the Tangent Gauss Map}
\author{Thomas Krämer, Vicente Cortés, Xiaoxiang Zhou}
\date{}

\begin{document}
\maketitle

The starting point of our discussion was the following question. Let $A$ be a simple abelian variety of dimension $n>2$, and let $C\subset A$ be an irreducible curve which generates $A$. Is the tangent Gauss map
\[
\gamma_C\colon C^{\mathrm{sm}}\longrightarrow \mathbb{P}^{n-1}
\]
generically injective unless the failure of injectivity is explained by a nontrivial symmetry of $C$ induced by an affine automorphism of $A$ whose linear part acts by a scalar?
\\

Our recent progress shows that this expectation is false. We construct a simple abelian threefold $A$ and a generating curve $C\subset A$ for which the tangent Gauss map has generic degree $3$, while no nontrivial affine automorphism of $A$ with scalar differential preserves $C$.
\\

We also discussed the main idea behind the construction. It starts from a special genus $4$ hyperelliptic curve arising from a tangential cover of an elliptic curve. We then deform this example to a non-hyperelliptic curve while preserving the Hodge-theoretic structure needed to obtain the corresponding abelian threefold. This leads to the desired counterexample and shows that non-injectivity of the tangent Gauss map can arise from the intrinsic geometry of the curve, rather than from symmetries of the ambient abelian variety.

\end{document}